\subsection{The complete exclusion of seventeen arrows}

\begin{theorem}[Proved: the unrestricted lower bound is eighteen arrows]
\label{CP2:A-global-lower-eighteen}
Every finite category whose classifying space is homotopy
equivalent to $\CP^2$ has at least eighteen arrows, counting
identities. In particular, the seventeen-arrow alternative
in Part~A.2 is false, and the least arrow number satisfies
\[
\boxed{18\leq n\leq22}.
\]
\end{theorem}
\begin{proof}
The complete bound through sixteen arrows is
Theorem~\ref{CP2:A-global-lower-seventeen}. It therefore
suffices to exclude exactly seventeen arrows. A category
of projective-plane type is connected; write $k$ for its
number of objects.

For $k=1$, the complete minimum-monoid classification
Theorem~\ref{A:monoid-minimum23-classification} excludes
all orders less than twenty-three. For $k=2$, use
Theorem~\ref{CP2:A-two-object-seventeen}; it includes
both support types and all object-retract reductions.

For $k=3$, cyclic support is excluded by
Theorem~\ref{CP2:A-strong-three-seventeen}. A connected
acyclic support is either a source or sink star, or a
full linear order: a directed two-step path forces its
third edge by composition. The star case is
Theorem~\ref{CP2:A-star-seventeen-complete}, and the
linear case is Theorem~\ref{CP2:A-linear-seventeen-complete}.

For $k=4$, use the cyclic-support theorem
\ref{CP2:A-strong-four-seventeen} and the complete
acyclic-support theorem \ref{CP2:A-four-seventeen-complete}.
Finally, if $k\geq5$, then $17\leq2k+7$ and $17\leq21$,
so the all-category excess theorem
\ref{CP2:A-high-all-eight} excludes the category. This
exhausts every finite category with seventeen arrows,
including empty cross hom-sets and nontrivial corners.

The independently checked construction of
Theorem~\ref{A:category-upper22} supplies the upper bound.
No inherited certificate stream absent from this run is
a premise of any finite gate used in this proof.
\end{proof}

\paragraph{Exact scope joins.}
The two-object gate has its separate root arithmetic and
receipt join in \path{outputs/two_object17_scope_independent.json}.
The star and linear gates cover all $924$ and $2240$
canonical cardinality profiles, respectively, with their
separate root joins in
\path{outputs/star17_scope_independent.json} and
\path{outputs/linear17_scope_independent.json}.

The final four-object scope has $34218$ canonical profiles.
Its original $33137$ structural cases are supplemented
by $56$ tree cases and $244$ actual singleton-comma
deletions, with $20$ overlapping cases, and the four
explicit final profiles. Thus its disjoint partition is
\[
34218=(33137+56+244-20+4)+797=33421+797.
\]
The required finite part contains $650840107$ weighted
corner-representative tables and $7267931$ literal
obstruction records. Of those records, $7247034$ have
a checked zero-square primitive and $20897$ have positive
third mod-two homology. There is no missing required
profile and no unexcluded record. These workload counts
are not counts of category isomorphism classes.

The separate root program
\path{code/check_four17_scope_join.py} imports no profile
enumerator or finalizer. It regenerates every canonical
cardinality row, reconstructs the three supplemental proof
sets, and checks all $797$ required receipts, input and
source hashes, compressed-stream digests and independent
literal-obstruction counts. Its complete receipt is
\path{outputs/four17_scope_independent.json}. Four historical
proof-redundant receipts whose failure-history filenames
changed are explicitly excluded from every adopted finite
count; their profiles are covered by the separately bound
tree proof. Their preserved records and subsequent history
recovery have no role in the required finite partition.

The unrestricted eighteen-arrow decision is the next
open case. No theorem in this subsection excludes an
eighteen-arrow category.
